% Proof note: O(g^4) phases at the <3>|<2> boundary. LaTeX fragment (REVTeX-compatible), not a full document.
% Numbers verified in exact rationals: scratchpad/domain_form.py (closed forms vs per_column for 1167 sequences),
% indep_rs.py (independent brute-force RS), e6_triples.py (O(g^6) range; see the last paragraph).

\paragraph*{Assumptions.}
(A1) Straight walls along $x$, $x$-uniform columns, parallel stacking for $r>0$, and $L_x\to\infty$. Moving a wall
by one row takes $L_x$ flips ($r=0$) or $2L_x$ flips ($r>0$), so different wall configurations do not mix at any
finite order, and each is treated by non-degenerate Rayleigh--Schr\"odinger theory at $T=0$.
(A2) Every domain has at least two sites. A one-site domain costs $O(1)$.
(A3) $\kappa=\tfrac12+kg^2+k_1g^4$ with $k=F/4$ and $F=6/[(2+r)(3+r)(5+r)]$, the $\langle3\rangle|\langle2\rangle$ line at
$O(g^2)$. On this line a domain of length $n\ge4$ costs $\tfrac{g^2}{2}[(n-2)F-F_4]>0$ at $O(g^2)$, with
$F_4=F-2/[(3+r)(4+r)(5+r)]$, whereas two- and three-site domains cost nothing. At $O(g^4)$ the competition is
therefore among sequences of 2s and 3s only. Their energy per site is $k_1(1-4\rho)+W$, with
$W=k\,e_2'+e_4$ evaluated at $\kappa=\tfrac12$.

\paragraph*{Locality.}
At $\kappa=\tfrac12$ and under (A2), the site $y\pm2$ lies either in the domain of $y$ or in the adjacent domain,
never further away. Hence the field $h_y$, the flip cost $D_y=2h_y$ and $c_y=s_y(s_{y+2}+s_{y-2})$ depend only on the
position of $y$ in its own domain (Table~\ref{tab:h}). The coupled pairs that enter $e_4$ are the $x$ bond and the
rung, which join sites of the same column position, and the $y$ and $y2$ bonds, whose two sites lie either in one
domain or on the two sides of a single wall. Therefore
\begin{equation}
W\,L=\sum_k\big[a(n_k)+b(n_k,n_{k+1})\big],
\end{equation}
where $L$ is the period in sites and $n_k$ the domain lengths. With
$P(D,D',B)=-4B(D+D')/[(DD')^2(D+D'-4B)]$, and $D_1,\dots,D_n$ the flip costs along a domain,
\begin{align}
a(n)&=\sum_{i=1}^{n}\Big[-\frac{2k\,c_i}{D_i^2}+\frac1{D_i^3}+P(D_i,D_i,1)+\tfrac12P(D_i,D_i,r)\Big]
      +\sum_{i=1}^{n-1}P(D_i,D_{i+1},1)+\sum_{i=1}^{n-2}P(D_i,D_{i+2},-\tfrac12),\nonumber\\
b(n,n')&=P(X_n,X_{n'},-1)+P(Y_n,X_{n'},\tfrac12)+P(X_n,Y_{n'},\tfrac12).
\end{align}
The flip costs along a domain are $2(3{+}r),2(3{+}r)$ for $n=2$; $2(2{+}r),2(5{+}r),2(2{+}r)$ for $n=3$; and
$2(2{+}r),2(4{+}r),2(3{+}r),\dots,2(4{+}r),2(2{+}r)$ for $n\ge4$. The values $c_i=-2,0,+2$ hold when both, one, or
none of the sites $i\pm2$ lie across a wall. $X_n$ and $Y_n$ are the flip costs of the first and second site from
the wall. Exact values:
$r=0$: $a(2)=\tfrac{7}{1080}$, $a(3)=-\tfrac{59}{2000}$, $b(2,2)=-\tfrac1{720}$, $b(2,3)=-\tfrac1{67200}$,
$b(3,3)=\tfrac3{400}$.
$r=1$: $a(2)=\tfrac1{768}$, $a(3)=-\tfrac{271}{60480}$, $b(2,2)=-\tfrac3{8960}$, $b(2,3)=\tfrac1{10368}$,
$b(3,3)=\tfrac5{3456}$.
How the energy is split between $a$ and $b$ is a convention; only the cycle costs below are invariant. These
formulas reproduce $W$ exactly for all 2/3 sequences with up to 10 domains, and for all sequences containing 4s
or 5s with up to 6 domains.

\paragraph*{Only $\langle2\rangle$, $\langle3\rangle$ and $\langle23\rangle$ can be stable.}
Since $k_1(1-4\rho)L=\sum_k k_1(n_k-4)$, the $O(g^4)$ energy per site is $\sum_k C(n_k,n_{k+1})/\sum_k n_k$ with
$C(n,n')=k_1(n-4)+a(n)+b(n,n')$. A periodic sequence is thus a closed walk on the complete directed graph, with
loops, on the vertices $\{2,3\}$. The edge $n\to n'$ has cost $C(n,n')$ and length $n$, and the energy per site is
total cost over total length. Every closed walk splits into simple cycles, and a ratio of sums with positive
denominators is at least the smallest ratio among the parts (the mediant inequality). Equality requires every part
to attain that smallest ratio. The simple cycles are $2\to2$, $3\to3$ and $2\to3\to2$, that is $\langle2\rangle$,
$\langle3\rangle$ and $\langle23\rangle$. So for every $k_1$ the minimum over all sequences, of any length, is
reached by one of these three. A sequence mixing different cycles can tie only on a phase boundary.

\paragraph*{Degenerate block mixtures.}
Let $N_{22}$, $N_{33}$ and $N_{23}=N_{32}$ count the consecutive domain pairs of a cyclic sequence. Then
$L=2N_{22}+3N_{33}+5N_{23}$, the number of walls is $N_{22}+N_{33}+2N_{23}$, and the total cost is
$N_{22}C_{2}+N_{33}C_{3}+N_{23}C_{23}$, with $C_2=C(2,2)$, $C_3=C(3,3)$ and $C_{23}=C(2,3)+C(3,2)$. Hence
$(\rho,W)$ is the convex combination of the points of $\langle2\rangle$, $\langle3\rangle$ and $\langle23\rangle$
with the site fractions $2N_{22}/L$, $3N_{33}/L$ and $5N_{23}/L$ as weights. Three consequences follow.
Sequences without ``33'' ($\langle223\rangle$, $\langle2223\rangle$, $\langle22323\rangle,\dots$) lie exactly on the
segment $\langle2\rangle$--$\langle23\rangle$. Sequences without ``22'' ($\langle233\rangle$, $\langle2333\rangle$,
$\langle23233\rangle,\dots$) lie exactly on the segment $\langle23\rangle$--$\langle3\rangle$. Sequences containing
both pairs lie strictly above the hull, because $\langle23\rangle$ lies strictly below the chord from
$\langle3\rangle$ to $\langle2\rangle$, by $-619/504000$ at $r=0$ and $-667/3628800$ at $r=1$.
With $W_{\langle2\rangle}=\tfrac{11}{4320}$, $W_{\langle3\rangle}=-\tfrac{11}{1500}$ and
$W_{\langle23\rangle}=-\tfrac{34849}{7560000}$ at $r=0$ (and $\tfrac{13}{26880}$, $-\tfrac{367}{362880}$,
$-\tfrac{2167}{3628800}$ at $r=1$), $\langle23\rangle$ is stable for
$k_1\in(\tfrac{20591}{2016000},\tfrac{6011}{336000})$ at $r=0$ and $(\tfrac{167}{107520},\tfrac{1961}{725760})$ at
$r=1$. Each window edge is itself a multiphase line at $O(g^4)$. On the $\langle23\rangle|\langle2\rangle$ edge all
33-free sequences are degenerate; on the $\langle3\rangle|\langle23\rangle$ edge all 22-free sequences are.

\paragraph*{Next order (range argument; see the numerical check below).}
At $O(g^6)$ the connected clusters have at most three sites, for example $\{y,y+2,y+4\}$, so they span at most one
intermediate domain, of length 2 or 3. The flip costs remain internal to their domains. The other $O(g^6)$ boundary
terms, $k_1e_2'$, $\tfrac{k^2}{2}e_2''$ and $k\,e_4'$, have range at most two domains. Hence
$W_6L=\sum_k t(n_{k-1},n_k,n_{k+1})$. Repeating the argument on the graph whose vertices are consecutive domain
pairs shows the following. On the $\langle23\rangle|\langle2\rangle$ edge (vertices 22, 23, 32) the only simple
cycles are $\langle2\rangle$, $\langle23\rangle$ and $\langle223\rangle$. On the $\langle3\rangle|\langle23\rangle$
edge they are $\langle3\rangle$, $\langle23\rangle$ and $\langle233\rangle$. So at $O(g^6)$ at most one new phase can
open in each gap: $\langle223\rangle$ or $\langle233\rangle$. No longer sequence can.
% Check (e6_triples.py, exact rationals): e6*L equals a sum over domain triples with zero residual for all 2/3
% sequences with up to 8 domains (71 sequences) at both r = 0 and r = 1, rank 5. The nearest-neighbour pair
% form fails (residual 4.5e-4 at r = 0, 2.5e-5 at r = 1). <2223> lies on the <2>-<223> segment of W6 to 1e-16
% (w6_indep.py).
